% 图契约：磁边界的法向与切向；图中的端口、几何、箭头和方程对应本题。
% 版面：沿用前三模块的几何图+推导分区；目标印宽17cm，简图至中等密度。
% 布局：左侧真实结构，右侧关键关系；真实电路边不作为架构流程箭头处理。
\documentclass[tikz,border=7pt]{standalone}
\usepackage{fontspec,xeCJK,amsmath,amssymb}
\IfFontExistsTF{Times New Roman}{\setmainfont{Times New Roman}}{\setmainfont{TeX Gyre Termes}}
\IfFontExistsTF{Heiti SC}{\setCJKmainfont[AutoFakeBold=1.5]{Heiti SC}}{\setCJKmainfont[AutoFakeBold=1.5]{FandolHei-Regular}}
\usetikzlibrary{arrows.meta,bending,calc,patterns,positioning,decorations.pathmorphing}
\definecolor{ink}{HTML}{243745}
\definecolor{blue}{HTML}{275DAB}
\definecolor{green}{HTML}{177C63}
\definecolor{red}{HTML}{BD4935}
\definecolor{orange}{HTML}{AA6C16}
\definecolor{purple}{HTML}{715393}
\tikzset{every node/.style={font=\fontsize{9.3}{12.5}\selectfont,text=ink},
arrow/.style={-{Stealth[length=5pt 1.5,width'=0pt 3]},shorten >=2pt,shorten <=1pt,line width=1pt},
arrow short/.style={-{Stealth[length=3pt 1.0,width'=0pt 2.5]},shorten >=1pt,shorten <=0pt,line width=.8pt},
axis/.style={arrow,draw=ink!60,line width=.6pt},
vec/.style={arrow,draw=blue},
wire/.style={draw=ink!75,line width=.8pt},
guide/.style={dashed,draw=ink!45,line width=.5pt},
card/.style={rounded corners=2pt,draw=blue!65,fill=blue!4,align=center,text width=6.5cm,inner xsep=9pt,inner ysep=7pt},
title/.style={font=\fontsize{10.5}{14}\selectfont\bfseries,text width=15cm,align=center},
label/.style={font=\fontsize{8.4}{11}\selectfont,align=center},
note/.style={label,text width=6.1cm}}
\begin{document}
\begin{tikzpicture}[x=1cm,y=1cm]
\path[use as bounding box] (-8,-3.5) rectangle(8,3.65);
\node[title] at (0,3.25) {M5-035\quad 磁边界的法向与切向};
\begin{scope}[shift={(-4.05,0)}]
\fill[blue!4](-2.8,-2)rectangle(2.8,0);\draw[ink](-2.8,0)--(2.8,0);\draw[guide](0,-2)--(0,2);\draw[vec,red](-.6,-1.03923)--(0,0);\draw[vec,red](0,0)--(2.4,1.03923);\draw[orange](0,-.7)arc(-90:-120:.7);\node[orange]at(-.45,-.95){$\alpha_1$};\draw[orange](0,.65)arc(90:23.413:.65);\node[orange]at(.75,.9){$\alpha_2$};\node[left]at(-2,-1){$\mu_1$};\node[left]at(-2,1){$\mu_2$};\node[right]at(.1,1.9){$\mathbf n_{12}$};\node[label]at(0,-2.5){$\mathbf K_f=0$；角度从法线量起};\end{scope}
\node[card,draw=blue!65,fill=blue!4] at (3.9,1.8) {$B_{1n}=B_{2n}$};
\node[card,draw=green!65,fill=green!4] at (3.9,0) {$\mathbf K_f=0:\ H_{1t}=H_{2t}$};
\node[card,draw=orange!65,fill=orange!4] at (3.9,-1.8) {$\mathbf B=\mu\mathbf H$};
\end{tikzpicture}\end{document}